\chapter{课程作业详解}
\label{app:homework}

本附录对应两份原始题单：第2章讲义PDF第64页，以及下载目录中的作业图片。两份题单第1--3题相同，但第4--5题不同，因此先讲共同题，再分别保留各自第4--5题，共七组解答。全章使用列向量、左乘矩阵和右手正旋转；行向量答案须将矩阵整体转置。

\section{共同第1题：世界、设备与规格化坐标}
\qs{题意}{世界坐标系、设备坐标系以及规格化坐标系三者的区别和联系是什么？}
\sol 世界坐标系是统一描述场景中各物体位置、方向及大小的全局参考，通常按工程尺度取单位；设备坐标系服务于实际显示或绘图设备，其单位、范围、原点与轴向依硬件而定；规格化设备坐标是二者之间与设备无关的中间表达，本课程取二维坐标范围$[0,1]^2$。

三者的联系是\textbf{由世界窗口到规格化坐标，再由规格化坐标映射到设备视口}。例如世界窗口$x\in[10,30]$、$y\in[20,60]$中的点$(15,30)$，规格化后为
\[
u=\frac{15-10}{30-10}=\frac14,\qquad
v=\frac{30-20}{60-20}=\frac14.
\]
若设备视口为$U\in[100,900]$、$V\in[50,650]$，且两个坐标系竖轴同向，则$(U,V)=(300,200)$；若设备竖轴向下而世界竖轴向上，则$V=650-600v=500$。
\nt{观察/相机坐标用于相对于观察者表达几何，设备坐标用于输出设备，不能把二者当成同一坐标系。窗口和视口宽高比不同且独立缩放时，图形可能发生非均匀变形。}
\sourcepage{第2章}{14--18、64}

\section{共同第2题：三角形比例变换}
\qs{题意}{已知$A(20,15)$、$B(20,40)$、$C(40,30)$。分别求：（1）$x$方向缩小一半，$y$方向增加一倍；（2）整个三角形放大为原来的1.5倍。}
\sol 按基本比例变换约定，以\textbf{坐标原点为缩放中心}；“增加一倍”表示变为原来的2倍，“放大为1.5倍”指各线性尺寸乘1.5。

（1）$s_x=1/2$、$s_y=2$，所以
\[
S_1=\begin{bmatrix}1/2&0&0\\0&2&0\\0&0&1\end{bmatrix},\qquad
\begin{bmatrix}x'\\y'\\1\end{bmatrix}=S_1\begin{bmatrix}x\\y\\1\end{bmatrix}.
\]
直接得到
\[
\boxed{A'(10,30),\quad B'(10,80),\quad C'(20,60).}
\]

（2）$s_x=s_y=1.5$，所以
\[
S_2=\begin{bmatrix}1.5&0&0\\0&1.5&0\\0&0&1\end{bmatrix},
\]
\[
\boxed{A''(30,22.5),\quad B''(30,60),\quad C''(60,45).}
\]
可用面积变化检验：二维比例变换的面积倍率是$|s_xs_y|$。原三角形面积为250，第一种仍为250，第二种为$250\times1.5^2=562.5$。第一种面积不变，但形状一般改变；第二种相似形状不变，面积不是只乘1.5。

\nt{若指定以某点$\boldsymbol q$为缩放中心，应改用$\boldsymbol p'=\boldsymbol q+S(\boldsymbol p-\boldsymbol q)$；不能把原点缩放结果直接用于其他中心。}
\sourcepage{第2章}{23、64}

\section{共同第3题：变换能否交换}
\qs{题意}{证明：（1）两个连续旋转；（2）两个连续平移；（3）两个连续比例变换；（4）比例系数相等时的旋转与比例变换，是否具有互换性。}
\sol 以下按本题的二维基本变换理解，旋转与比例均关于同一坐标原点，比例轴为同一组$x,y$轴。

\subsection*{（1）同中心二维旋转可以交换}
由三角函数和角公式，
\[
R(\beta)R(\alpha)=
\begin{bmatrix}
\cos(\alpha+\beta)&-\sin(\alpha+\beta)&0\\
\sin(\alpha+\beta)&\cos(\alpha+\beta)&0\\
0&0&1
\end{bmatrix}
=R(\alpha)R(\beta).
\]
故复合角为$\alpha+\beta$，顺序无关。\textbf{若旋转中心不同，或为三维不同轴旋转，一般不能交换。}

\subsection*{（2）两次平移可以交换}
\[
T(a_2,b_2)T(a_1,b_1)=
\begin{bmatrix}1&0&a_1+a_2\\0&1&b_1+b_2\\0&0&1\end{bmatrix}
=T(a_1,b_1)T(a_2,b_2).
\]
向量相加满足交换律，因此总平移量与先后次序无关。

\subsection*{（3）沿同一坐标轴的比例变换可以交换}
\[
S(c_2,d_2)S(c_1,d_1)=
\operatorname{diag}(c_2c_1,d_2d_1,1)
=S(c_1,d_1)S(c_2,d_2).
\]
这是对角矩阵乘法的结果。若缩放中心或主轴不同，不能直接使用该结论。

\subsection*{（4）等比例缩放与同原点旋转可以交换}
取$S(k,k)=\operatorname{diag}(k,k,1)$，令$c=\cos\theta,s=\sin\theta$，则
\[
S(k,k)R(\theta)=
\begin{bmatrix}kc&-ks&0\\ks&kc&0\\0&0&1\end{bmatrix}
=R(\theta)S(k,k).
\]
一般非均匀比例变换不成立。例如对$(1,0)$，取$S(2,1)$与$R(90^\circ)$，先缩放后旋转得$(0,2)$，反序得$(0,1)$。
\nt{证明“可交换”需说明对象、中心和轴的条件。不能把二维同中心旋转的结论推广为“三维任意旋转都可交换”。}
\sourcepage{第2章}{30--33、64}

\section{讲义第4题：裁剪的概念与分类}
\qs{题意}{论述图形裁剪的概念，并按被裁剪对象分类。}
\sol 图形裁剪是按照给定的显示窗口或可视区域，确定并保留图形内部部分、舍弃外部部分的处理过程。二维通常相对窗口裁剪，三维则相对可视体裁剪。

按被裁剪对象，可分为\textbf{点裁剪、直线段裁剪、多边形裁剪、曲线裁剪、曲面裁剪}等。点只需进行区域判定；线段可能需要计算与边界的交点；多边形除求交外，还要组织有序的输出顶点，维持边界结构。图形在窗口内但被另一物体遮挡，是消隐处理的问题，不能只靠裁剪解决。
\sourcepage{第2章}{44--49、64}

\section{讲义第5题：Cohen--Sutherland编码与求交}
\qs{题意}{（1）给出九个区域的四位编码；（2）窗口左下角$(0,0)$、右上角$(5,5)$，判断$A_1(3,6)$至$B_1(-1,-1)$，以及$A_2(-1,7)$至$B_2(6,6)$两条线段与窗口的关系，并给出过程。}
\sol 编码从高位到低位为上、下、右、左。边界点视为窗口内，因此九区为
\[
\begin{array}{c|c|c}
1001&1000&1010\\\hline
0001&0000&0010\\\hline
0101&0100&0110
\end{array}
\qquad
\begin{array}{l}
\text{列分界：}x=0,\ x=5,\\
\text{行分界：}y=5,\ y=0.
\end{array}
\]

\subsection*{线段$A_1B_1$：部分可见}
$A_1$在上方，编码$1000$；$B_1$在左下方，编码$0101$。按位与为$0000$，按位或非零，需继续求交。

先取上边界$y=5$。原线段参数式为
\[
\boldsymbol p(t)=(3,6)+t(-4,-7),\qquad0\leq t\leq1.
\]
由$6-7t=5$，得$t=1/7$，交点
\[
I_{\rm top}=\left(\frac{17}{7},5\right),
\]
编码为$0000$。用该点替换$A_1$。

再处理$B_1$的左侧编码，与$x=0$求交：$3-4t=0$，得$t=3/4$，
\[
I_{\rm left}=\left(0,\frac34\right),
\]
编码也为$0000$。此时两个新端点均在窗口内，接受线段
\[
\boxed{\left(\frac{17}{7},5\right)\longleftrightarrow\left(0,\frac34\right).}
\]
也可由各边界不等式得到$1/7\leq t\leq3/4$进行核对。

\subsection*{线段$A_2B_2$：全部不可见}
$A_2$编码$1001$，$B_2$编码$1010$，所以
\[
1001\mathbin{\&}1010=1000\ne0000.
\]
两端点都位于上边界外侧，整条线段的$y$坐标在6至7之间，故不可能进入$y\leq5$的窗口，直接舍弃。
\sourcepage{第2章}{48、64}

\section{图片第4题：绕任意轴旋转后再绕$z$轴旋转}
\qs{题意}{三角形顶点为$P_1=(0,1,0)$、$P_2=(2,1,0)$、$P_3=(2,3,0)$，先绕$P_1P_3$旋转$90^\circ$，再绕$z$轴旋转$45^\circ$。写出基本变换矩阵和最终复合变换矩阵，计算各点最终坐标。}
\sol 原题未指定有向轴与正旋转方向。这里约定第一步沿$\overrightarrow{P_1P_3}$为轴正向，第二步沿固定世界$+z$轴，两个角度均取\textbf{右手正向}。若取反向轴，第一步相当于改为$-90^\circ$，结果将不同。

\subsection*{1. 平移并把任意轴对准$z$轴}
轴上点与单位方向为
\[
\boldsymbol q=P_1=(0,1,0)^{\mathsf T},\qquad
\boldsymbol u=\frac{P_3-P_1}{\|P_3-P_1\|}
=\frac{1}{\sqrt2}(1,1,0)^{\mathsf T}.
\]
先用$T(0,-1,0)$把轴移过原点，再绕$z$轴$-45^\circ$使$\boldsymbol u$对准$x$轴，最后绕$y$轴$-90^\circ$使其对准$+z$轴。记
\[
Q=R_y(-90^\circ)R_z(-45^\circ),\qquad Q\boldsymbol u=(0,0,1)^{\mathsf T}.
\]
本题用到的基本齐次矩阵是平移$T(0,\pm1,0)$与$\widehat R_z(\pm45^\circ)$、$\widehat R_y(\pm90^\circ)$、$\widehat R_z(90^\circ)$；其中帽号表示给$3\times3$旋转矩阵补齐次末行末列。

令$a=1/\sqrt2$，具体的旋转块为
\[
R_z(\pm45^\circ)=\begin{bmatrix}a&\mp a&0\\\pm a&a&0\\0&0&1\end{bmatrix},\quad
R_y(\pm90^\circ)=\begin{bmatrix}0&0&\pm1\\0&1&0\\\mp1&0&0\end{bmatrix},
\]
\[
R_z(90^\circ)=\begin{bmatrix}0&-1&0\\1&0&0\\0&0&1\end{bmatrix},\qquad
T(0,\pm1,0)=\begin{bmatrix}1&0&0&0\\0&1&0&\pm1\\0&0&1&0\\0&0&0&1\end{bmatrix}.
\]

\subsection*{2. 绕轴旋转，再撤销对准和平移}
列向量依次经历的基本变换给出
\begin{align*}
H_{\rm axis}={}&T(0,1,0)\widehat R_z(45^\circ)\widehat R_y(90^\circ)
\widehat R_z(90^\circ)\\[-1mm]
&\qquad\cdot\widehat R_y(-90^\circ)\widehat R_z(-45^\circ)T(0,-1,0).
\end{align*}
其旋转块也可由Rodrigues公式检验：
\[
R=\boldsymbol u\boldsymbol u^{\mathsf T}+[\boldsymbol u]_\times
=\begin{bmatrix}1/2&1/2&a\\1/2&1/2&-a\\-a&a&0\end{bmatrix}.
\]
因此
\[
H_{\rm axis}=\begin{bmatrix}R&\boldsymbol q-R\boldsymbol q\\\boldsymbol0^{\mathsf T}&1\end{bmatrix}
=\begin{bmatrix}
1/2&1/2&a&-1/2\\
1/2&1/2&-a&1/2\\
-a&a&0&-a\\
0&0&0&1
\end{bmatrix}.
\]
第一步结果为
\[
P_1'=(0,1,0),\qquad P_2'=(1,2,-\sqrt2),\qquad P_3'=(2,3,0).
\]
$P_1,P_3$位于轴上，应保持不动，这是检查任意轴旋转的直接方法。

\subsection*{3. 绕固定$z$轴旋转$45^\circ$}
第二步左乘$\widehat R_z(45^\circ)$，得到总矩阵
\[
\boxed{H=\widehat R_z(45^\circ)H_{\rm axis}
=\begin{bmatrix}
0&0&1&-a\\
a&a&0&0\\
-a&a&0&-a\\
0&0&0&1
\end{bmatrix}.}
\]
所以对任意点，
\[
x''=z-a,\qquad y''=a(x+y),\qquad z''=a(-x+y-1).
\]
代入三点：
\[
\boxed{P_1''=\left(-\frac1{\sqrt2},\frac1{\sqrt2},0\right),\quad
P_2''=\left(-\frac1{\sqrt2},\frac3{\sqrt2},-\sqrt2\right),}
\]
\[
\boxed{P_3''=\left(-\frac1{\sqrt2},\frac5{\sqrt2},0\right).}
\]
检验三边长度：$\|P_2''-P_1''\|=2$、$\|P_3''-P_2''\|=2$、$\|P_3''-P_1''\|=2\sqrt2$，与原三角形一致。第二步后轴上两点也随世界$z$轴旋转，因此不能要求它们最终仍停在原坐标。

\nt{若第一步按$\overrightarrow{P_3P_1}$正向，$P_2'$的$z$分量变为$+\sqrt2$，而第二步后$P_2''$的$z$也为$+\sqrt2$。写清轴向后，才有唯一的有符号答案。}
{\footnotesize\color{mygr}题目来源：CADCAM目录中的作业图片，第4题。}

\section{图片第5题：旋转矩阵的逆与性质}
\qs{题意}{旋转矩阵$R$的逆矩阵是什么？旋转矩阵还有哪些性质？}
\sol 对右手系中的纯旋转，$R\in\mathbb R^{3\times3}$满足
\[
\boxed{R^{-1}=R^{\mathsf T},\qquad R^{\mathsf T}R=RR^{\mathsf T}=I,\qquad\det R=+1.}
\]
原因是旋转把一组正交单位基变为另一组同手性的正交单位基。矩阵各列（各行）互相正交且为单位向量，故$R^{\mathsf T}R=I$；保持右手性排除了$\det R=-1$的反射。

因此旋转具有下列性质：
\begin{enumerate}
\item \textbf{保内积、长度、距离与夹角：}
$(R\boldsymbol a)^{\mathsf T}(R\boldsymbol b)=\boldsymbol a^{\mathsf T}\boldsymbol b$，因而$\|R\boldsymbol a\|=\|\boldsymbol a\|$。
\item \textbf{保方向性与体积：}$\det R=1$，不改变有向体积与坐标系手性；且$R(\boldsymbol a\times\boldsymbol b)=(R\boldsymbol a)\times(R\boldsymbol b)$。
\item \textbf{同轴逆旋转：}$R(\boldsymbol u,\theta)^{-1}=R(\boldsymbol u,-\theta)$。
\item \textbf{复合仍是旋转：}若$R_1,R_2$为旋转矩阵，$R_2R_1$仍正交且行列式为1，但通常$R_2R_1\ne R_1R_2$。
\item \textbf{旋转轴不变：}三维纯旋转存在单位轴向$\boldsymbol u$满足$R\boldsymbol u=\boldsymbol u$；旋转角可由$\operatorname{tr}R=1+2\cos\theta$检验。
\end{enumerate}

若讨论的是带平移的齐次刚体变换，不能把整个矩阵的逆误写为转置：
\[
\boxed{\begin{bmatrix}R&\boldsymbol t\\\boldsymbol0^{\mathsf T}&1\end{bmatrix}^{-1}
=\begin{bmatrix}R^{\mathsf T}&-R^{\mathsf T}\boldsymbol t\\\boldsymbol0^{\mathsf T}&1\end{bmatrix}.}
\]
逆变换先撤销平移，再反向旋转，所以平移块是$-R^{\mathsf T}\boldsymbol t$，并非一般意义上的$-\boldsymbol t$。
{\footnotesize\color{mygr}题目来源：CADCAM目录中的作业图片，第5题；基本旋转矩阵见第2章讲义PDF第26、37页。}
